%------------------------------------------------------------------------------
% GENERATED FILE -- DO NOT EDIT DIRECTLY.
% Source of truth: data/profile.yaml
% Regenerate with: python3 scripts/cv.py build
% Template derived from Jake Gutierrez / sb2nov (MIT License).
%------------------------------------------------------------------------------

\documentclass[letterpaper,10pt]{article}

\usepackage[T1]{fontenc}
\usepackage[utf8]{inputenc}
\usepackage{lmodern}
\usepackage{latexsym}
\usepackage[empty]{fullpage}
\usepackage{titlesec}
\usepackage{marvosym}
\usepackage[usenames,dvipsnames]{color}
\usepackage{verbatim}
\usepackage{enumitem}
\usepackage[hidelinks]{hyperref}
\usepackage{fancyhdr}
\usepackage[french]{babel}
\frenchsetup{StandardLists=true}
\usepackage{tabularx}
\usepackage{fontawesome5}
\input{glyphtounicode}

\pagestyle{fancy}
\fancyhf{}
\fancyfoot{}
\renewcommand{\headrulewidth}{0pt}
\renewcommand{\footrulewidth}{0pt}

% Margins
\addtolength{\oddsidemargin}{-0.5in}
\addtolength{\evensidemargin}{-0.5in}
\addtolength{\textwidth}{1.0in}
\addtolength{\topmargin}{-0.55in}
\addtolength{\textheight}{1.1in}

\urlstyle{same}

\raggedbottom
\raggedright
\setlength{\tabcolsep}{0in}

% Section formatting
\titleformat{\section}{
  \vspace{-5pt}\scshape\raggedright\large
}{}{0em}{}[\color{black}\titlerule \vspace{-5pt}]

% Machine-readable / ATS-parsable PDF
\pdfgentounicode=1

% Custom commands
\newcommand{\resumeItem}[1]{
  \item\small{
    {#1\vspace{-2pt}}
  }
}

\newcommand{\resumeSubheading}[4]{
  \vspace{-2pt}\item
    \begin{tabular*}{0.97\textwidth}[t]{l@{\extracolsep{\fill}}r}
      \textbf{#1} & #2 \\
      \textit{\small#3} & \textit{\small #4} \\
    \end{tabular*}\vspace{-7pt}
}

\newcommand{\resumeProjectHeading}[2]{
    \item
    \begin{tabular*}{0.97\textwidth}{l@{\extracolsep{\fill}}r}
      \small#1 & #2 \\
    \end{tabular*}\vspace{-7pt}
}

\renewcommand\labelitemii{$\vcenter{\hbox{\tiny$\bullet$}}$}

\newcommand{\resumeSubHeadingListStart}{\begin{itemize}[leftmargin=0.15in, label={}]}
\newcommand{\resumeSubHeadingListEnd}{\end{itemize}}
\newcommand{\resumeItemListStart}{\begin{itemize}[itemsep=-1pt]}
\newcommand{\resumeItemListEnd}{\end{itemize}\vspace{-5pt}}

\begin{document}

%----------HEADING----------
\begin{center}
    \textbf{\Huge \scshape Ilias Laoukili} \\ \vspace{2pt}
    \href{https://www.linkedin.com/in/ilias-laoukili}{\color{MidnightBlue}\faLinkedin\ linkedin.com/in/ilias-laoukili} $|$
    \href{mailto:ilias.laoukili@proton.me}{\color{MidnightBlue}\faEnvelope\ ilias.laoukili@proton.me} $|$
    \href{https://github.com/ilias-laoukili}{\color{MidnightBlue}\faGithub\ github.com/ilias-laoukili} $|$
    \href{https://ilias.laoukili.com}{\color{MidnightBlue}\faGlobe\ ilias.laoukili.com} \\
    \small Québec, QC, Canada
\end{center}

\section{Profil professionnel}
\small{
    Élève-ingénieur en IA et cybersécurité et étudiant à la maîtrise recherche (M.Sc.) à l'Université Laval, actuellement conseiller en intelligence artificielle au ministère de la Cybersécurité et du Numérique (MCN). Expérience pratique en systèmes RAG, GNN et bibliothèques Rust/Python haute performance (recherche au LAMA, stage R\&D à Nanyang Polytechnic, Singapour).
}
\vspace{-4pt}

\section{Formation}
\resumeSubHeadingListStart
  \resumeSubheading
    {Université Laval}{Québec, QC, Canada}
    {Maîtrise en informatique avec mémoire (M.Sc. recherche) -- Double diplôme}{Automne 2026 -- Présent}
    \resumeItemListStart
      \resumeItem{Maîtrise recherche en informatique et IA, sous la direction du Prof. Talal Halabi, en complément du cursus d'ingénieur.}
    \resumeItemListEnd

  \resumeSubheading
    {ESIEE Paris}{Noisy-le-Grand, France}
    {Diplôme d'ingénieur en Intelligence Artificielle et Cybersécurité (Grande École)}{2024 -- 2027}
    \resumeItemListStart
      \resumeItem{Diplôme d'ingénieur reconnu par l'État (CTI). Cours : apprentissage automatique, vision par ordinateur, science des données, cybersécurité.}
    \resumeItemListEnd

  \resumeSubheading
    {Lycée Gustave Monod}{Enghien-les-Bains, France}
    {Classes Préparatoires aux Grandes Écoles (CPGE)}{2021 -- 2024}
    \resumeItemListStart
      \resumeItem{Préparation intensive en mathématiques et physique (MPSI/PCSI/PC) aux concours des grandes écoles d'ingénieurs.}
    \resumeItemListEnd

\resumeSubHeadingListEnd

\section{Expérience professionnelle}
\resumeSubHeadingListStart
  \resumeSubheading
    {Conseiller en intelligence artificielle (Emploi étudiant)}{Nov. 2026 -- Présent}
    {Ministère de la Cybersécurité et du Numérique (MCN) -- Bureau de la transformation de l'IA}{Québec, QC, Canada}
    \resumeItemListStart
      \resumeItem{Contribution à la capacité ministérielle de réalisation et d'industrialisation de services numériques IA sécurisés, opérationnels et réutilisables.}
      \resumeItem{Soutien à la documentation d'architecture, la préparation des données, le prototypage, les essais et l'intégration de solutions IA avec des spécialistes (architecture, données, sécurité, MLOps).}
    \resumeItemListEnd

  \resumeSubheading
    {Stagiaire ingénieur R\&D}{Mai 2026 -- Août 2026}
    {Nanyang Polytechnic (NYP)}{Singapour}
    \resumeItemListStart
      \resumeItem{Développement d'un système IA pour la simulation énergétique de bâtiments durables, avec conception d'un framework \textbf{RAG} rendant \textbf{EnergyPlus} interrogeable en langage naturel.}
      \resumeItem{Collaboration avec les équipes de recherche pour optimiser les simulations d'efficacité énergétique par des algorithmes d'IA appliquée.}
    \resumeItemListEnd

  \resumeSubheading
    {Assistant de recherche (Tremplin Recherche)}{Nov. 2025 -- Mai 2026}
    {Université Gustave Eiffel / LAMA}{Champs-sur-Marne, France}
    \resumeItemListStart
      \resumeItem{Étude des compromis entre performance et efficacité de calcul des réseaux de neurones sur graphes (GNN) via la sparsification de graphes.}
      \resumeItem{Analyse de la topologie des graphes et des schémas de sparsification afin de réduire le coût de calcul sans perte de précision.}
    \resumeItemListEnd

  \resumeSubheading
    {Tuteur académique et professeur particulier}{Sept. 2021 -- Mai 2026}
    {Complétude, Parkours, ESIEE Paris}{Île-de-France, France}
    \resumeItemListStart
      \resumeItem{Accompagnement de plus de 45 élèves en mathématiques et informatique, avec une amélioration moyenne des résultats de 30\%.}
    \resumeItemListEnd

\resumeSubHeadingListEnd

\section{Projets}
\resumeSubHeadingListStart
  \resumeProjectHeading
    {\textbf{\href{https://github.com/ilias-laoukili/plekton}{\color{MidnightBlue}plekton -- Bibliothèque Rust/Python de graphes synthétiques et sparsification}} $|$ \emph{Rust, PyO3, Python, GNN}}{2026}
    \resumeItemListStart
      \resumeItem{Conception d'une bibliothèque en Rust pur avec bindings Python zero-copy pour la génération de graphes synthétiques, l'ajustement de modèles à blocs stochastiques, les métriques structurelles et la sparsification spectrale.}
      \resumeItem{Pensée pour le benchmarking reproductible et rapide de GNN avec intégration directe aux pipelines ML Python.}
    \resumeItemListEnd

  \resumeProjectHeading
    {\textbf{\href{https://github.com/ilias-laoukili/sentim-app}{\color{MidnightBlue}Sentim -- Analyse d'émotions audio par IA}} $|$ \emph{Python, PyTorch, Librosa, Streamlit}}{Nov. 2025}
    \resumeItemListStart
      \resumeItem{Conception d'une plateforme complète de classification de 8 émotions vocales (RAVDESS), combinant extraction DSP (MFCC, Chroma, Tonnetz) et modèles CNN/ResNet sur spectrogrammes 2D surpassant la baseline Random Forest.}
      \resumeItem{Déploiement d'un tableau de bord Streamlit interactif pour la visualisation acoustique en temps réel et la manipulation prosodique.}
    \resumeItemListEnd

\resumeSubHeadingListEnd

\section{Compétences techniques}
\begin{itemize}[leftmargin=0.15in, label={}]
  \small{\item{
    \textbf{IA \& Langages clés}{: Python, Rust, SQL, PyTorch, PyTorch Geometric, Scikit-learn, GNN, RAG/LLM, DSP (Librosa)} \\
    \textbf{Outils \& DevOps}{: Git/GitHub, Docker, CI/CD, PyO3, Streamlit, Hydra, Linux, Wireshark, MongoDB, EnergyPlus} \\
    \textbf{Certifications}{: MongoDB Associate Developer, Harvard CS50's \textit{Introduction to AI with Python}} \\
    \textbf{Langues}{: Français (Langue maternelle), Allemand (C1 -- Abibac), Anglais (Compétence professionnelle -- TOEIC 935/990)}
  }}
\end{itemize}

\section{Engagement associatif}
\resumeSubHeadingListStart
  \resumeSubheading
    {Bénévole}{2023 -- Juin 2026}
    {Croix-Rouge française}{Puteaux, France}
    \resumeItemListStart
      \resumeItem{Contribution à l'aide sociale et à la logistique de terrain auprès de plus de 200 familles.}
    \resumeItemListEnd

\resumeSubHeadingListEnd

\end{document}
